\scanpage{111}
3. Decomposition of Ramification Index
\mspage{3.1}

The main purpose of this section is to prove Theorem 3.8, which provides
a decomposition of the ramification index into the sum of certain integers.

Given a local $R$-algebra $S$ over a local ring $R$, we have defined
$e(S/R)=\dfrac{e_S(mS)}{e_R(m)}$ and called it the multiplicity of the local $R$-algebra
$S$, where $m$ is the maximal ideal of $R$. As we have seen, the basic object
in the theory of ramification index is the multiplicity of local algebras.
This number $e(S/R)$ is not in general an integer ([9]) and a natural
question is to ask under what circumstances this number is an integer.
We already know two cases:

1) If $R$ is a regular local ring, then $e(S/R)$ is an integer since
$e_R(\underline m)=1$ where $\underline m$ is the maximal ideal of $R$.

2) If $R\subset S$ is a local algebra of normal local rings such that $\overline S$ is
separably algebraic over $\overline R$, then $e(S/R)$ is equal to $r(S/R)$ by Theorem 2.7
and hence is an integer.

The following propositions enable us to relax the hypothesis in the
above two cases.

Lemma 3.1: If $S$ is a local $R$-algebra over a local ring $R$ such that
$\operatorname{w.hd}_R S<\infty$, then $\widehat S$ is a finite local $\widehat R$-algebra such that $\operatorname{hd}_{\widehat R}\widehat S<\infty$.

\uline{Proof}: By the definition of a local algebra, $S=A_{\underline p}$ for a finite
$R$-algebra $A$ where $\underline p$ is a maximal ideal in $A$. It follows easily that
$\widehat S=(A\otimes_R\widehat R)_{\widehat{\underline p}}=(S\otimes_R\widehat R)_{\underline p'}$ where $\widehat{\underline p}=\underline p\otimes_R\widehat R$ and $\underline p'$ is the maximal ideal
in $S\otimes_R\widehat R$ generated by $\widehat{\underline p}$. Therefore $\widehat S$ is flat over $S\otimes_R\widehat R$ and consequently
$\operatorname{w.hd}_{\widehat R}\widehat S\leq\operatorname{w.hd}_{\widehat R}S\otimes_R\widehat R\leq\operatorname{w.hd}_R S$.
Thus if $\operatorname{w.hd}_R S<\infty$, then $\operatorname{w.hd}_{\widehat R}\widehat S<\infty$.
However $\widehat S$ is finitely generated as $\widehat R$-module since $\overline S$ is a finitely generated
$R$-module, and consequently $\operatorname{w.hd}_{\widehat R}\widehat S<\infty$ implies $\operatorname{hd}_{\widehat R}\widehat S<\infty$.

Proposition 3.2: Let $S$ be a local $R$-algebra over a local ring $R$ such
that $\operatorname{w.hd}_R S<\infty$. Then $e(S/R)$ is an integer, equal to $\sum_i(-1)^iL_S(:\operatorname{Tor}_i^R(S\ \overline R))$
where $L_S(\ )$ denotes the length as $S$-module.
\mspage{110}
